\documentclass[11pt]{article}
\usepackage[margin=25mm]{geometry}
\usepackage{fontspec}
\setmainfont{Latin Modern Roman}
\usepackage{amsmath,amssymb,mathtools}
\usepackage{unicode-math}
\setmathfont{Latin Modern Math}
\usepackage{xurl,hyperref}
\hypersetup{colorlinks=true,urlcolor=blue}
\setcounter{secnumdepth}{0}
\setlength{\parindent}{0pt}
\setlength{\parskip}{5pt}
\setlength{\emergencystretch}{3em}
\newcommand{\sref}[2]{\hyperlink{source:#1:#2}{[S#2]}}
\newcommand{\cataloguescope}[1]{\paragraph{Recorded target.}#1}
\newcommand{\R}{\mathbb R}
\newcommand{\C}{\mathbb C}
\newcommand{\N}{\mathbb N}
\newcommand{\Z}{\mathbb Z}
\newcommand{\Q}{\mathbb Q}
\newcommand{\F}{\mathbb F}
\newcommand{\eps}{\varepsilon}
\begin{document}
% UnsolvedMath ID 30003733; source catalogue code OWR-15993-001
\setcounter{section}{30003732}
\section{Low-Parameter Hyperelastic Models of Brain Tissue}\label{30003733}
{\small UnsolvedMath ID 30003733 · Mathematical Physics}
\subsection{Definitions and mathematical statement}

\paragraph{Shared notation.}$\mathbb N=\{1,2,\ldots\}$; the standard number fields have their usual meanings. All additional conventions are stated locally.


\paragraph{Definitions and conventions.}
A hyperelastic material has strain energy $W(F)$ depending on deformation gradient $F$; stress follows by differentiating this energy. Objectivity and isotropy impose invariance under appropriate rotations of $F$. Shear and axial compression or tension probe different combinations of its invariants. A low-parameter family uses few calibratable constitutive constants; validation requires comparison with independent multiaxial data, not merely interpolation of one loading curve.

\paragraph{Target retained from the canonical source catalogue.}
\paragraph{Catalogue wording retained for traceability.}
Determine a low-parameter family of isotropic hyperelastic strain-energy functions for brain tissue that captures its observed response under combined shear and axial compression or tension, remains suitable for small-to-moderate in-vivo strains, and can be calibrated and validated on multiaxial data.

\paragraph{Source scope and qualification.}
This is an experimentally constrained constitutive-model question. The report\textquotesingle s observed shear/compression/tension trends and the small-to-moderate in-vivo strain regime are part of the target, not a request for a universal energy independent of tissue or data. ID 30003734 is a duplicate extraction.
\subsection{Short English statement}
Find a small, testable hyperelastic model that captures the unusual combined-load response of brain tissue.
\subsection{Research attempt}
\paragraph{Model, date and investigated scope.}
OpenAI GPT-6 Astra (configured model \texttt{gpt-6-astra}), ultra reasoning; generated 11 October 2026. Derived shear nonidentifiability and complementary uniaxial parameter separation for an incompressible two-invariant candidate, without data validation.

\paragraph{Sources and literature check.}
Read the supplied definition, catalogue wording, source qualifications and reference entries for 30003733. The historical source/context and inherited status are retained; no fresh literature or novelty certification is claimed. No external theorem is used without proof in this attempt; explicit model assumptions are distinguished from facts about the source model.

\paragraph{An identifiability calculation.} I test the familiar two-parameter incompressible isotropic candidate \(W=c_1(I_1-3)+c_2(I_2-3)\), where \(I_1=\operatorname{tr}(F^TF)\), \(I_2=\operatorname{tr}((F^TF)^{-1})\) when \(\det F=1\). This is a candidate family, not a claimed fit to brain data. Along simple shear
\[
 F=\begin{pmatrix}1&k&0\\0&1&0\\0&0&1\end{pmatrix},
 \qquad I_1=I_2=3+k^2,
\]
so \(W=(c_1+c_2)k^2\) and the work-conjugate shear stress is \(2(c_1+c_2)k\). Shear observations alone cannot identify the two constants: replacing \((c_1,c_2)\) by \((c_1+a,c_2-a)\) leaves the entire shear curve unchanged.
\paragraph{A complementary load.} For incompressible uniaxial stretch \(F=\operatorname{diag}(\lambda,\lambda^{-1/2},\lambda^{-1/2})\), \(\lambda>0\),
\[
 \frac{dW}{d\lambda}=2c_1(\lambda-\lambda^{-2})+2c_2(1-\lambda^{-3}).
\]
The ratio of these two coefficient functions is \(\lambda\) whenever \(\lambda\ne1\). Thus the shear sum plus one nontrivial uniaxial measurement determine both constants algebraically, whereas the stress at \(\lambda=1\) supplies no information. Near \(1\) the coefficient functions nearly coincide, explaining an elementary conditioning problem.
\paragraph{Remaining target.} This is a calibration diagnostic, not validation of observed compression/shear/tension trends, tissue heterogeneity or the in-vivo strain range in [S2]. No experimental data are supplied here and none are fabricated. Establishing that any low-parameter candidate fits independent combined-load observations while retaining admissibility is the first gap. The calculation demonstrates why a successful single-mode fit would be insufficient and does not claim new constitutive modelling progress.

\paragraph{Verification and adversarial audit.}
Computed both invariants directly, checked lambda=1 degeneracy and the coefficient ratio, and distinguished work-conjugate path stress from arbitrary stress components. This is a GPT-6 Astra Ultra self-audit, not independent human verification. Prompt A drafting followed by prompt B scope/dependency/boundary self-audit. Computed both invariants directly, checked lambda=1 degeneracy and the coefficient ratio, and distinguished work-conjugate path stress from arbitrary stress components. The remaining gap is: Fit and validate a mechanically admissible low-parameter energy on independent combined-load brain-tissue measurements. Unresolved status and zero original-target completion are retained.

Coordinator recomputed the simple-shear invariants and uniaxial derivatives, including the lambda=1 degeneracy and coefficient ratio lambda. No data-fit claim is supported. This additional check was model-assisted, not independent human verification.

\paragraph{Tools, process and contributions.}
Python 3 and local file-reading/writing tools for manuscript serialization and structural checks; mathematical calculations were symbolic reasoning, with no numerical experiment or theorem prover. Contribution prompt A was followed by prompt B. GPT-6 Astra with ultra reasoning drafted the research and performed its stated self-audit; any additional coordinator contribution is identified in the verification and audit. A GPT-6-based Codex coordinator assembled the manuscript and checked repository integration; its exact variant and reasoning configuration were not exposed. No human mathematical contribution, independent human proof check or formal proof certification is supplied. Complete derivations and the consulted source references are included here; private reasoning transcripts are not published.

\paragraph{Final status and first remaining gap.}
\textbf{UNRESOLVED}. Fit and validate a mechanically admissible low-parameter energy on independent combined-load brain-tissue measurements.

\section{Completion Estimate}
\noindent\textbf{COMPLETION: 0\%}
The proved calculation is an elementary auxiliary result or a test of a restricted strategy; no new portion of the full source target is established.

\subsection{Sources}
\begin{itemize}
\item[\hypertarget{source:30003733:1}{S1}] ULAM.AI and the UnsolvedMath Contributors, UnsolvedMath, record 30003733 (OWR-15993-001). Catalogue curation under CC BY 4.0; definitions and English exposition added locally.
\url{https://huggingface.co/datasets/ulamai/UnsolvedMath}.
\item[\hypertarget{source:30003733:2}{S2}] The Mathematics of Mechanobiology and Cell Signaling (2018), Oberwolfach Reports, DOI 10.4171/owr/2018/8, pp. 433--506. Primary report accessed 2026-10-11; relevant PDF pages 12.
\url{https://doi.org/10.4171/owr/2018/8}.
\item[\hypertarget{source:30003733:3}{S3}] original OWR report (catalogue-listed context; not independently consulted in this preparation).
\url{https://doi.org/10.4171/OWR/2018/8}.
\item[\hypertarget{source:30003733:4}{S4}] Mihai et al., hyperelastic brain models (catalogue-listed context; not independently consulted in this preparation).
\url{https://doi.org/10.1016/j.jmps.2017.05.015}.
\item[\hypertarget{source:30003733:5}{S5}] Filla et al., heterogeneous brain constitutive model (catalogue-listed context; not independently consulted in this preparation).
\url{https://doi.org/10.1016/j.jmbbm.2023.106271}.
\end{itemize}
\label{end:30003733}
\end{document}
